\section{模型评价与推广}

\subsection{模型优点}

\begin{enumerate}[label=\arabic*)]
  \item \textbf{守恒口径统一。} 从干固体和水分两条连续方程出发，区分干基含水率与单位体积浓度；问题四在材料坐标中不重复添加压缩项，并用累计收支验证。
  \item \textbf{四问衔接清楚。} 问题一、二、四从相同原始初态独立计算，问题三沿问题二轨迹继续；固定域、耦合、事件和收缩的新增因素可逐层解释。
  \item \textbf{判停方式稳健。} 终点由全域最大含水率决定，并以非中心湿峰反例排除“只看中心”的程序漏洞；临界等号与严格达标整数秒分开报告。
  \item \textbf{验证来源相互独立。} 解析基线、独立有限元、网格与时间收敛、累计守恒和输入留出分别检查方程、离散、实现与数据处理。
\end{enumerate}

\subsection{模型局限}

\begin{enumerate}[label=\arabic*)]
  \item 缺少药材内部温湿度和总质量实测，无法校准$C_{\mathrm{eq}}=G(Y,T)$、物性误差或给出统计置信区间。
  \item 主热方程没有蒸发潜热、水分携带焓、辐射和形变功，属于有效显热模型，而非完整多相能量模型。
  \item 一维主模型假定端面封闭且轴向均匀；二维对照只覆盖“端面与侧面同样暴露”这一情景，并非未知边界的严格上下界。
  \item 问题四半径由观测直接指定，模型预测给定收缩轨迹下的温湿场，尚未建立含水率与收缩之间的闭合关系；经验密度作为有效热学参数使用。
  \item 4 h后的外部环境使用均值延拓，虽已作情景分析，但缺少后续观测验证。
\end{enumerate}

\subsection{改进与推广}

若获得实验数据，可优先测量不同温湿环境下的吸附/解吸等温线与内部温湿剖面，用于辨识$C_{\mathrm{eq}}$、$D(C,T)$及其不确定性；进一步将潜热和水分焓通量纳入能量方程，并通过称重数据校准干质量与体积关系。若端面实际暴露，应使用经绝对误差验证的二维轴对称模型；若存在轴向各向异性或非均匀收缩，则扩展到二维各向异性或三维移动网格模型。

这一框架可用于根茎、果蔬等近圆柱多孔材料。更换物性、平衡含水率关系和几何观测后，守恒方程、材料坐标及全域事件的求解方式仍可沿用。用于工艺控制前，还需依据校准结果和测量误差设置停机裕度。

\section{结论}

本文以干基含水率守恒为主线，完成了圆柱药材由预热、耦合干燥到径向收缩条件下的统一建模，主要结论如下：

\begin{enumerate}[label=\arabic*)]
  \item 问题一的热、水分扩散长度在30 min时分别约为$1.743$ cm和$0.298$ cm。数值结果与尺度判断一致：温度影响已进入内部，显著失水仍集中在表层，因此温度趋稳不能代表水分均匀。
  \item 问题二中，3 h时中心---表面温差仅$0.1169\ ^\circ\mathrm C$，含水率差仍为$0.7581$。冻结扩散率的温度因子会明显减少前三小时失水，说明所选经验模型中的温度依赖对早期过程影响显著。
  \item 问题三以全域最大含水率定义事件，得到等号临界时间$57.4740$ h，首个严格达标整数秒为$206907$ s。平均值判停会提前约$21.45$ h，说明最后湿点决定了实际终点。
  \item 问题四在材料坐标中得到临界时间$51.0920$ h。在统一外生半径轨迹下，四组条件对照表明，相对问题三缩短的$6.3820$ h同时包含几何、物性和交互影响。两问对扩散率尺度的局部灵敏度均约为$-0.89$。
\end{enumerate}

解析基线、网格收敛、独立有限元和累计守恒均支持上述数值结果。宽情景计算进一步表明，当长期平衡含水率接近$0.15$时，终点会显著延后，甚至超出半径观测时域。实际应用应优先校准$D(C,T)$、平衡含水率和表面传质关系，再结合安全裕度确定停机时间。
